\begin{table}[htbp]
\centering
\caption{Predicted probability of usable height data by residence and survey period}
\label{tab:pp_residence}
\begin{threeparttable}

\small
\renewcommand{\arraystretch}{1.15}
\setlength{\tabcolsep}{25pt}

\begin{tabularx}{\textwidth}{>{\raggedright\arraybackslash}l r r r}
\toprule
\addlinespace[1pt]
\textbf{Survey period} & \textbf{Urban} & \textbf{Rural} &
\makecell[rt]{\textbf{Rural--Urban}\\\textbf{difference}} \\
\addlinespace[1pt]
\midrule
2011--2016 & 0.901 (0.896, 0.907) & 0.922 (0.920, 0.924) & 0.021 \\
2017--2022 & 0.905 (0.900, 0.910) & 0.930 (0.928, 0.932) & 0.025 \\
\bottomrule
\end{tabularx}

\begin{tablenotes}[flushleft]
\footnotesize
\item Notes: Predicted probabilities are marginal estimates from a survey-weighted logistic regression including a residence $\times$ survey period interaction. Values are predicted probabilities with 95\% confidence intervals. Joint Wald test for interaction: \(p=0.0844\).
\end{tablenotes}

\end{threeparttable}
\end{table}